% Report for the 2-D transient heat conduction project (AIAA conference format).
%
% Needs new-aiaa.cls and new-aiaa.bst in this folder. They are not part of TeX Live:
% open the AIAA template in the Overleaf gallery, "Download as zip", and copy both files here.
%
% Build from this folder:   latexmk -pdf report.tex      Clean up:   latexmk -c
% Figures come from ../figs, written by `.venv/bin/python plot.py` from the project root.

\documentclass[conf]{new-aiaa}
%\documentclass[journal]{new-aiaa} for journal papers
\usepackage[utf8]{inputenc}

\usepackage{graphicx}
\usepackage{amsmath}
\usepackage{siunitx}
\usepackage{longtable,tabularx}
\setlength\LTleft{0pt}

\graphicspath{{../figs/}}

\title{Transient Two-Dimensional Heat Conduction in a Rectangular Bar:\\ Explicit FTCS and ADI Solutions}


\author{Samuel R. Denis\footnote{Graduate Student, Department of Aerospace Engineering, Wichita State University.}}
\affil{Department of Aerospace Engineering, Wichita State University, Wichita, Kansas, 67260}

\begin{document}

\maketitle

\begin{abstract}
The transient temperature field in a long copper bar with fixed wall temperatures is 
computed with the explicit forward-time, centered-space (FTCS) scheme and 
the alternating direction implicit (ADI) scheme.
The steady state is compared with the analytical methods on a grid of 31 x 41.
With a time step of 0.2s both schemes reach steady-state, which is defined as the total variation in
temperature per step being less than \SI{0.01}{\degreeCelsius} 
defined as CONSS after about 410 s of simulated time.
For both methods the max error is \SI{0.2874}{\degreeCelsius} near the hot corners where the wall discontinuity exists, 
root-mean-square error of \SI{0.0260}{\degreeCelsius}, mean error of \SI{-0.0011}{\degreeCelsius}.
With a timestep of 1.0 s, FTCS diverges, as its stability limit predicts, ADI converges in 496 steps. 
The steady-state wall heat flux satisfies the overall energy balance of the system within \SI{0.1}{\percent}.
\end{abstract}

\section*{Nomenclature}

{\renewcommand\arraystretch{1.0}
\noindent\begin{longtable*}{@{}l @{\quad=\quad} l@{}}
$L$, $W$ & bar length in $x$ and width in $y$, \si{\metre} \\
$T$ & temperature, \si{\degreeCelsius} \\
$T_0$ & initial interior temperature, \si{\degreeCelsius} \\
$T_1$, $T_2$, $T_3$, $T_4$ & bottom, left, top, and right wall temperatures, \si{\degreeCelsius} \\
$\alpha$ & thermal diffusivity, \si{\metre\squared\per\second} \\
$k$ & thermal conductivity, \si{\watt\per\metre\per\degreeCelsius} \\
$t$ & time, \si{\second} \\
$\Delta t$ & time step, \si{\second} \\
$\Delta x$, $\Delta y$ & grid spacing in $x$ and $y$, \si{\metre} \\
$i$, $j$ & grid indices in $x$ and $y$ \\
$n$ & time level \\
$d_x$, $d_y$ & FTCS diffusion numbers, $\alpha\Delta t/\Delta x^2$ and $\alpha\Delta t/\Delta y^2$ \\
$d_1$, $d_2$ & ADI half-step coefficients, $\alpha\Delta t/(2\Delta x^2)$ and $\alpha\Delta t/(2\Delta y^2)$ \\
$G$ & amplification factor \\
$\mathrm{TV}$ & total variation, sum of $|T^{n+1}_{i,j} - T^{n}_{i,j}|$ over interior nodes, \si{\degreeCelsius} \\
$\mathrm{CONSS}$ & steady-state tolerance on TV, \si{\degreeCelsius} \\
$m$ & series term index \\
$q_{\mathrm{in}}$ & heat flux into the bar, \si{\watt\per\metre\squared}
\end{longtable*}}

\section{Introduction}
This report presents solutions to Problems 3.5 through 3.8 of Hoffmann and Chiang~\cite{hoffmann2000}.
Problem 3.5 describes transient heat conduction in a long rectangular copper bar, initially at a
uniform temperature, the four sides are held at constant temperatures. The transient solution is
computed with the explicit forward-time, centered-space (FTCS) scheme, the alternating direction,
implicit (ADI) scheme, or the fractional step method (FSM), and the steady state is compared with the
analytical series solution. The steady-state heat flux is evaluated along the sides $y = 0$, $y = W$, and
$x = 0$. The ADI method was selected over FSM. Finally, the accuracy, stability, and efficiency of the three schemes are compared. 

\section{Problem Statement}
% TODO: geometry, material, initial and boundary conditions.
The geometry shown in Fig.~\ref{fig:geometry} defines the axis system as x-right and y-up
with the length and width of the bar defined as L and W and respectively. 
The figure defines the indices i and j in base one 
whereas in the solution provided all indices are given in base 0 notation with the ends of the bar
defined as IM and JM respectively. 
\begin{figure}[hbt!]
\centering
\includegraphics[width=.5\textwidth]{Figure-P3_5}
\caption{Geometry, axes, and boundary conditions of the bar (after Ref.~\cite{hoffmann2000}).}
\label{fig:geometry}
\end{figure}
Fig.~\ref{fig:geometry}

The temperature in the bar satisfies the two-dimensional heat conduction equation
\begin{equation}
\label{eq:heat}
\frac{\partial T}{\partial t} = \alpha \left( \frac{\partial^2 T}{\partial x^2} + \frac{\partial^2 T}{\partial y^2} \right)
\end{equation}
on $0 \le x \le L$, $0 \le y \le W$, with $T = T_0$ initially and fixed wall temperatures
$T(x,0) = T_1$, $T(0,y) = T_2$, $T(x,W) = T_3$, and $T(L,y) = T_4$.
It is assumed that the solution has reached steady state 
if the total variation in temperature from one time point to the next is less than CONSS. CONSS is given as 
\SI{0.01}{\degreeCelsius} in the problem statement. 
\section{Numerical Methods}

\subsection{Explicit FTCS Scheme}
% TODO: derivation and stability limit.
\begin{equation}
\label{eq:ftcs}
T^{n+1}_{i,j} = T^{n}_{i,j}
 + d_x \left( T^{n}_{i+1,j} - 2T^{n}_{i,j} + T^{n}_{i-1,j} \right)
 + d_y \left( T^{n}_{i,j+1} - 2T^{n}_{i,j} + T^{n}_{i,j-1} \right)
\end{equation}
The scheme is stable for $d_x + d_y \le 1/2$.

\subsection{Alternating Direction Implicit Scheme}
% TODO: derivation; why each half step is tridiagonal; wall values moved into the right-hand side.
\begin{align}
\label{eq:adi-x}
-d_1 T^{*}_{i-1,j} + (1 + 2d_1) T^{*}_{i,j} - d_1 T^{*}_{i+1,j}
 &= d_2 T^{n}_{i,j-1} + (1 - 2d_2) T^{n}_{i,j} + d_2 T^{n}_{i,j+1} \\
\label{eq:adi-y}
-d_2 T^{n+1}_{i,j-1} + (1 + 2d_2) T^{n+1}_{i,j} - d_2 T^{n+1}_{i,j+1}
 &= d_1 T^{*}_{i-1,j} + (1 - 2d_1) T^{*}_{i,j} + d_1 T^{*}_{i+1,j}
\end{align}

\subsection{Thomas Algorithm}
% TODO: forward elimination and back substitution for each tridiagonal line; O(n) cost.
The ADI method as discussed splits the pentadiagonal system into a pair of tridiagonal systems of equations. The naive approach to solving a system of 
equations would be to employ gaussian elimination which has a big O complexity of $O(n^3)$. To efficiently 
solve the special case of a tridiagonal system Eq.~\eqref{eq:tridiag} the Thomas Algorithm [Thomas 1949] was developed. 
Hoffmann and Chiang~\cite{hoffmann2000} demonstrates the same algorithm in Appendix B. under a different name.
This approach takes advantage of the zero symmetry about the diagonal to solve the system in $O(n)$ complexity.
The notation below matches the code provided. 
% ---------- Tridiagonal system ----------
\begin{equation}
  a_k\,x_{k-1} + b_k\,x_k + c_k\,x_{k+1} = d_k,
  \qquad k = 0, 1, \dots, n-1,
  \qquad a_0 = c_{n-1} = 0
  \label{eq:tridiag}
\end{equation}

% ---------- Matrix form ----------
\begin{equation}
  \begin{bmatrix}
    b_0    & c_0    &        &         &         \\
    a_1    & b_1    & c_1    &         &         \\
           & \ddots & \ddots & \ddots  &         \\
           &        & a_{n-2}& b_{n-2} & c_{n-2} \\
           &        &        & a_{n-1} & b_{n-1}
  \end{bmatrix}
  \begin{bmatrix} x_0 \\ x_1 \\ \vdots \\ x_{n-2} \\ x_{n-1} \end{bmatrix}
  =
  \begin{bmatrix} d_0 \\ d_1 \\ \vdots \\ d_{n-2} \\ d_{n-1} \end{bmatrix}
  \label{eq:tridiag-matrix}
\end{equation}

Where the $a_k$,$b_k$,$c_k$ is known and $d_k$ is unknown.
For each time step the inner nodes of a row form a system of equations Eq.~\eqref{eq:tridiag-matrix}.
The algorithm begins with a linear forwards sweep handling the first equation which does not have $a_k$ 
by using equations Eq.~\eqref{eq:thomas-first}. Then the rest of the rows forward sweep is solved of the form Eq.~\eqref{eq:thomas-forward}.
% ---------- Forward sweep ----------
\begin{equation}
  s_0 = \frac{c_0}{b_0}, \qquad
  x'_0 = \frac{d_0}{b_0}
  \label{eq:thomas-first}
\end{equation}
\begin{equation}
  \left.
  \begin{aligned}
    \delta_k &= b_k - a_k\,s_{k-1} \\
    s_k      &= \frac{c_k}{\delta_k}, \qquad k \le n-2 \\
    x'_k     &= \frac{d_k - a_k\,x'_{k-1}}{\delta_k}
  \end{aligned}
  \right\}
  \qquad k = 1, 2, \dots, n-1
  \label{eq:thomas-forward}
\end{equation}
The loop completes over the inner nodes and then back substitutes following Eq.~\eqref{eq:thomas-back}.
% ---------- Back substitution ----------
\begin{equation}
  x_{n-1} = x'_{n-1}, \qquad
  x_k = x'_k - s_k\,x_{k+1}, \qquad k = n-2, n-3, \dots, 0
  \label{eq:thomas-back}
\end{equation}

The algorithm returns with the solved rows unknown $d_k$.

\subsection{Analytical Steady-State Solution}
% TODO: derivation by separation of variables (valid for T_2 = T_4 = 0); truncation criterion.
The analytical steady state solution provided by the problem statement can be simplified for the case of $T_2 = T_4 = 0$ as Eq.~\eqref{eq:analytical}.
The solution below is continuous and was evaluated at the same grid of nodes as the solvers for consistency of comparison.
Numerical stability of the hyperbolic sin had to be taken into account due to double precision floating point overflow.
The identity $\sinh a / \sinh b = e^{a-b} (1 - e^{-2a}) / (1 - e^{-2b})$, for $0 \le a \le b$, enables the stable calculation of Eq.~\eqref{eq:analytical}.


\begin{equation}
\label{eq:analytical}
T(x,y) = \sum_{m\ \mathrm{odd}} \frac{4}{m\pi} \sin\!\left(\frac{m\pi x}{L}\right)
\frac{ T_1 \sinh\!\left[ m\pi (W - y)/L \right] + T_3 \sinh\!\left( m\pi y/L \right) }{ \sinh\!\left( m\pi W/L \right) }
\end{equation}

\subsection{Convergence and Wall Heat Flux}
% TODO: steady-state criterion TV < CONSS; second-order one-sided wall gradient and sign convention.
\begin{equation}
\label{eq:flux}
q_{\mathrm{in}} = -k\,\nabla T \cdot \hat{n}_{\mathrm{in}}, \qquad
\left.\frac{\partial T}{\partial y}\right|_{y=0} \approx \frac{-3T_{i,0} + 4T_{i,1} - T_{i,2}}{2\Delta y}
\end{equation}

\section{Results}
% Order: the analytical baseline first, then FTCS, then ADI, then the FTCS/ADI comparison.

\subsection{Analytical Baseline}
% TODO: the steady-state field every numerical result is compared against; series truncation.

\begin{figure}[hbt!]
\centering
\includegraphics[width=.5\textwidth]{fig_analytical}
\caption{Analytical steady-state temperature field.}
\label{fig:analytical}
\end{figure}

\subsection{FTCS}

\subsubsection{Transient Solution}
% TODO: Part I discussion.

\begin{figure}[hbt!]
\centering
\includegraphics[width=\textwidth]{fig_ftcs_dt0.2_snapshots}
\caption{FTCS temperature field at $t = \SI{10}{\second}$, $t = \SI{40}{\second}$, and steady state, $\Delta t = \SI{0.2}{\second}$.}
\label{fig:ftcs-snapshots}
\end{figure}

\begin{figure}[hbt!]
\centering
\includegraphics[width=\textwidth]{fig_ftcs_dt0.2_profiles}
\caption{FTCS temperature versus $y$ at $x = 0$ to \SI{0.30}{\metre} in steps of \SI{0.05}{\metre}, $\Delta t = \SI{0.2}{\second}$.}
\label{fig:ftcs-profiles}
\end{figure}

\subsubsection{Stability}
% TODO: d_x + d_y for each time step; growth rate versus the amplification factor.

\begin{figure}[hbt!]
\centering
\includegraphics[width=\textwidth]{fig_ftcs_dt1_instability}
\caption{FTCS with $\Delta t = \SI{1.0}{\second}$: growth of the probe temperature compared with the amplification factor $|G|^n$, and the first 40 steps.}
\label{fig:ftcs-instability}
\end{figure}

\subsubsection{Temperature Histories and Convergence}
% TODO: Part III discussion; decay rate versus the slowest mode.

\begin{figure}[hbt!]
\centering
\includegraphics[width=.5\textwidth]{fig_ftcs_dt0.2_probes}
\caption{FTCS temperature histories at the three probe locations, $\Delta t = \SI{0.2}{\second}$.}
\label{fig:ftcs-probes}
\end{figure}

\begin{figure}[hbt!]
\centering
\includegraphics[width=.5\textwidth]{fig_ftcs_dt0.2_convergence}
\caption{FTCS total variation versus time with the steady-state tolerance and the slowest-mode decay rate.}
\label{fig:ftcs-convergence}
\end{figure}

\subsubsection{Steady State Compared with the Analytical Solution}
% TODO: Part II discussion; error metrics; corner singularity.

\begin{figure}[hbt!]
\centering
\includegraphics[width=\textwidth]{fig_ftcs_dt0.2_steady_comparison}
\caption{Analytical and FTCS steady-state fields and their difference.}
\label{fig:ftcs-steady-comparison}
\end{figure}

\begin{figure}[hbt!]
\centering
\includegraphics[width=.5\textwidth]{fig_ftcs_dt0.2_centerline}
\caption{Steady-state temperature along the centerline $x = \SI{0.15}{\metre}$, with the FTCS error.}
\label{fig:ftcs-centerline}
\end{figure}

\begin{figure}[hbt!]
\centering
\includegraphics[width=.5\textwidth]{fig_ftcs_dt0.2_near_wall_error}
\caption{FTCS steady-state error along the first interior row, $y = \SI{0.01}{\metre}$.}
\label{fig:ftcs-near-wall-error}
\end{figure}

\subsection{ADI}

\subsubsection{Transient Solution}
% TODO: ADI transient results; agreement with FTCS at t = 10 s and 40 s.

\begin{figure}[hbt!]
\centering
\includegraphics[width=\textwidth]{fig_adi_dt0.2_snapshots}
\caption{ADI temperature field at $t = \SI{10}{\second}$, $t = \SI{40}{\second}$, and steady state, $\Delta t = \SI{0.2}{\second}$.}
\label{fig:adi-snapshots}
\end{figure}

\begin{figure}[hbt!]
\centering
\includegraphics[width=\textwidth]{fig_adi_dt0.2_profiles}
\caption{ADI temperature versus $y$ at $x = 0$ to \SI{0.30}{\metre} in steps of \SI{0.05}{\metre}, $\Delta t = \SI{0.2}{\second}$.}
\label{fig:adi-profiles}
\end{figure}

\subsubsection{Temperature Histories and Convergence}
% TODO: ADI histories; convergence at a time step where FTCS is unstable.

\begin{figure}[hbt!]
\centering
\includegraphics[width=.5\textwidth]{fig_adi_dt0.2_probes}
\caption{ADI temperature histories at the three probe locations, $\Delta t = \SI{0.2}{\second}$.}
\label{fig:adi-probes}
\end{figure}

\begin{figure}[hbt!]
\centering
\includegraphics[width=.5\textwidth]{fig_adi_dt0.2_convergence}
\caption{ADI total variation versus time with the steady-state tolerance and the slowest-mode decay rate.}
\label{fig:adi-convergence}
\end{figure}

\subsubsection{Steady State Compared with the Analytical Solution}
% TODO: ADI error metrics; same grid error as FTCS.

\begin{figure}[hbt!]
\centering
\includegraphics[width=\textwidth]{fig_adi_dt0.2_steady_comparison}
\caption{Analytical and ADI steady-state fields and their difference.}
\label{fig:adi-steady-comparison}
\end{figure}

\begin{figure}[hbt!]
\centering
\includegraphics[width=.5\textwidth]{fig_adi_dt0.2_centerline}
\caption{Steady-state temperature along the centerline $x = \SI{0.15}{\metre}$, with the ADI error.}
\label{fig:adi-centerline}
\end{figure}

\subsection{Comparison of FTCS and ADI}
% TODO: 3.8: accuracy at t = 10 s and 40 s, steps and cost to steady state,
% ADI at a time step where FTCS is unstable. The overlay figures are not generated yet.

\subsection{Wall Heat Flux}
% TODO: placement undecided; move this subsection wherever it reads best.
% Part IV discussion; sign convention; energy balance; corner spikes grow with refinement.

\begin{figure}[hbt!]
\centering
\includegraphics[width=\textwidth]{fig_dt0.2_heat_flux}
\caption{Steady-state heat flux into the bar along the bottom, top, and left walls (top row) and the difference from the analytical field (bottom row). Corners are excluded; the flux there is unbounded.}
\label{fig:heat-flux}
\end{figure}

\section{Conclusion}
% TODO: main findings; do not repeat the abstract.

\section*{Acknowledgments}
Claude (Anthropic), used through Claude Code, assisted with this work. It supported the writing of the Python
post-processing and plotting scripts, wrote portions of the output and file-handling code in the Zig solver,
and reviewed and supported the debugging of the author's implementations of the FTCS, ADI, Thomas, and analytical
solutions. It also supported the initial \LaTeX{} structure, validation of LaTeX equations, figure captions of this
report and educated on the AIAA format, which the author revised. The author
verified all results, wrote the remaining text, and is responsible for all of the content provided.

\bibliography{references}

\end{document}
